\documentclass{report}

%%%%%%%%%%%%%%%%%%%%%%%%%%%%%%%%%
% PACKAGE IMPORTS
%%%%%%%%%%%%%%%%%%%%%%%%%%%%%%%%%


\usepackage[tmargin=2cm,rmargin=1in,lmargin=1in,margin=0.85in,bmargin=2cm,footskip=.2in]{geometry}
\usepackage{amsmath,amsfonts,amsthm,amssymb,mathtools}
\usepackage[varbb]{newpxmath}
\usepackage{xfrac}
\usepackage[makeroom]{cancel}
\usepackage{mathtools}
\usepackage{bookmark}
\usepackage{enumitem}
\usepackage{hyperref,theoremref}
\hypersetup{
	pdftitle={Assignment},
	colorlinks=true, linkcolor=doc!90,
	bookmarksnumbered=true,
	bookmarksopen=true
}
\usepackage[most,many,breakable]{tcolorbox}
\usepackage{xcolor}
\usepackage{varwidth}
\usepackage{varwidth}
\usepackage{etoolbox}
%\usepackage{authblk}
\usepackage{nameref}
\usepackage{multicol,array}
\usepackage{tikz-cd}
\usepackage[ruled,vlined,linesnumbered]{algorithm2e}
\usepackage{comment} % enables the use of multi-line comments (\ifx \fi) 
\usepackage{import}
\usepackage{xifthen}
\usepackage{pdfpages}
\usepackage{transparent}
\usepackage[utf8]{inputenc}
\usepackage{fontawesome}


\newcommand\mycommfont[1]{\footnotesize\ttfamily\textcolor{blue}{#1}}
\SetCommentSty{mycommfont}
\newcommand{\incfig}[1]{%
    \def\svgwidth{\columnwidth}
    \import{./figures/}{#1.pdf_tex}
}

\usepackage{tikzsymbols}
\renewcommand{\qedsymbol}{\faWhatsapp}


%\usepackage{import}
%\usepackage{xifthen}
%\usepackage{pdfpages}
%\usepackage{transparent}


%%%%%%%%%%%%%%%%%%%%%%%%%%%%%%
% SELF MADE COLORS
%%%%%%%%%%%%%%%%%%%%%%%%%%%%%%



\definecolor{myg}{RGB}{56, 140, 70}
\definecolor{myb}{RGB}{45, 111, 177}
\definecolor{myr}{RGB}{199, 68, 64}
\definecolor{mytheorembg}{HTML}{F2F2F9}
\definecolor{mytheoremfr}{HTML}{00007B}
\definecolor{mylenmabg}{HTML}{FFFAF8}
\definecolor{mylenmafr}{HTML}{983b0f}
\definecolor{mypropbg}{HTML}{f2fbfc}
\definecolor{mypropfr}{HTML}{191971}
\definecolor{myexamplebg}{HTML}{F2FBF8}
\definecolor{myexamplefr}{HTML}{88D6D1}
\definecolor{myexampleti}{HTML}{2A7F7F}
\definecolor{mydefinitbg}{HTML}{E5E5FF}
\definecolor{mydefinitfr}{HTML}{3F3FA3}
\definecolor{notesgreen}{RGB}{0,162,0}
\definecolor{myp}{RGB}{197, 92, 212}
\definecolor{mygr}{HTML}{2C3338}
\definecolor{myred}{RGB}{127,0,0}
\definecolor{myyellow}{RGB}{169,121,69}
\definecolor{myexercisebg}{HTML}{F2FBF8}
\definecolor{myexercisefg}{HTML}{88D6D1}


%%%%%%%%%%%%%%%%%%%%%%%%%%%%
% TCOLORBOX SETUPS
%%%%%%%%%%%%%%%%%%%%%%%%%%%%

\setlength{\parindent}{1cm}
%================================
% THEOREM BOX
%================================

\tcbuselibrary{theorems,skins,hooks}
\newtcbtheorem[number within=section]{Theorem}{Theorem}
{%
	enhanced,
	breakable,
	colback = mytheorembg,
	frame hidden,
	boxrule = 0sp,
	borderline west = {2pt}{0pt}{mytheoremfr},
	sharp corners,
	detach title,
	before upper = \tcbtitle\par\smallskip,
	coltitle = mytheoremfr,
	fonttitle = \bfseries\sffamily,
	description font = \mdseries,
	separator sign none,
	segmentation style={solid, mytheoremfr},
}
{th}

\tcbuselibrary{theorems,skins,hooks}
\newtcbtheorem[number within=chapter]{theorem}{Theorem}
{%
	enhanced,
	breakable,
	colback = mytheorembg,
	frame hidden,
	boxrule = 0sp,
	borderline west = {2pt}{0pt}{mytheoremfr},
	sharp corners,
	detach title,
	before upper = \tcbtitle\par\smallskip,
	coltitle = mytheoremfr,
	fonttitle = \bfseries\sffamily,
	description font = \mdseries,
	separator sign none,
	segmentation style={solid, mytheoremfr},
}
{th}


\tcbuselibrary{theorems,skins,hooks}
\newtcolorbox{Theoremcon}
{%
	enhanced
	,breakable
	,colback = mytheorembg
	,frame hidden
	,boxrule = 0sp
	,borderline west = {2pt}{0pt}{mytheoremfr}
	,sharp corners
	,description font = \mdseries
	,separator sign none
}

%================================
% Corollery
%================================
\tcbuselibrary{theorems,skins,hooks}
\newtcbtheorem[number within=section]{Corollary}{Corollary}
{%
	enhanced
	,breakable
	,colback = myp!10
	,frame hidden
	,boxrule = 0sp
	,borderline west = {2pt}{0pt}{myp!85!black}
	,sharp corners
	,detach title
	,before upper = \tcbtitle\par\smallskip
	,coltitle = myp!85!black
	,fonttitle = \bfseries\sffamily
	,description font = \mdseries
	,separator sign none
	,segmentation style={solid, myp!85!black}
}
{th}
\tcbuselibrary{theorems,skins,hooks}
\newtcbtheorem[number within=chapter]{corollary}{Corollary}
{%
	enhanced
	,breakable
	,colback = myp!10
	,frame hidden
	,boxrule = 0sp
	,borderline west = {2pt}{0pt}{myp!85!black}
	,sharp corners
	,detach title
	,before upper = \tcbtitle\par\smallskip
	,coltitle = myp!85!black
	,fonttitle = \bfseries\sffamily
	,description font = \mdseries
	,separator sign none
	,segmentation style={solid, myp!85!black}
}
{th}


%================================
% LENMA
%================================

\tcbuselibrary{theorems,skins,hooks}
\newtcbtheorem[number within=section]{Lenma}{Lenma}
{%
	enhanced,
	breakable,
	colback = mylenmabg,
	frame hidden,
	boxrule = 0sp,
	borderline west = {2pt}{0pt}{mylenmafr},
	sharp corners,
	detach title,
	before upper = \tcbtitle\par\smallskip,
	coltitle = mylenmafr,
	fonttitle = \bfseries\sffamily,
	description font = \mdseries,
	separator sign none,
	segmentation style={solid, mylenmafr},
}
{th}

\tcbuselibrary{theorems,skins,hooks}
\newtcbtheorem[number within=chapter]{lenma}{Lenma}
{%
	enhanced,
	breakable,
	colback = mylenmabg,
	frame hidden,
	boxrule = 0sp,
	borderline west = {2pt}{0pt}{mylenmafr},
	sharp corners,
	detach title,
	before upper = \tcbtitle\par\smallskip,
	coltitle = mylenmafr,
	fonttitle = \bfseries\sffamily,
	description font = \mdseries,
	separator sign none,
	segmentation style={solid, mylenmafr},
}
{th}


%================================
% PROPOSITION
%================================

\tcbuselibrary{theorems,skins,hooks}
\newtcbtheorem[number within=section]{Prop}{Proposition}
{%
	enhanced,
	breakable,
	colback = mypropbg,
	frame hidden,
	boxrule = 0sp,
	borderline west = {2pt}{0pt}{mypropfr},
	sharp corners,
	detach title,
	before upper = \tcbtitle\par\smallskip,
	coltitle = mypropfr,
	fonttitle = \bfseries\sffamily,
	description font = \mdseries,
	separator sign none,
	segmentation style={solid, mypropfr},
}
{th}

\tcbuselibrary{theorems,skins,hooks}
\newtcbtheorem[number within=chapter]{prop}{Proposition}
{%
	enhanced,
	breakable,
	colback = mypropbg,
	frame hidden,
	boxrule = 0sp,
	borderline west = {2pt}{0pt}{mypropfr},
	sharp corners,
	detach title,
	before upper = \tcbtitle\par\smallskip,
	coltitle = mypropfr,
	fonttitle = \bfseries\sffamily,
	description font = \mdseries,
	separator sign none,
	segmentation style={solid, mypropfr},
}
{th}


%================================
% CLAIM
%================================

\tcbuselibrary{theorems,skins,hooks}
\newtcbtheorem[number within=section]{claim}{Claim}
{%
	enhanced
	,breakable
	,colback = myg!10
	,frame hidden
	,boxrule = 0sp
	,borderline west = {2pt}{0pt}{myg}
	,sharp corners
	,detach title
	,before upper = \tcbtitle\par\smallskip
	,coltitle = myg!85!black
	,fonttitle = \bfseries\sffamily
	,description font = \mdseries
	,separator sign none
	,segmentation style={solid, myg!85!black}
}
{th}



%================================
% Exercise
%================================

\tcbuselibrary{theorems,skins,hooks}
\newtcbtheorem[number within=section]{Exercise}{Exercise}
{%
	enhanced,
	breakable,
	colback = myexercisebg,
	frame hidden,
	boxrule = 0sp,
	borderline west = {2pt}{0pt}{myexercisefg},
	sharp corners,
	detach title,
	before upper = \tcbtitle\par\smallskip,
	coltitle = myexercisefg,
	fonttitle = \bfseries\sffamily,
	description font = \mdseries,
	separator sign none,
	segmentation style={solid, myexercisefg},
}
{th}

\tcbuselibrary{theorems,skins,hooks}
\newtcbtheorem[number within=chapter]{exercise}{Exercise}
{%
	enhanced,
	breakable,
	colback = myexercisebg,
	frame hidden,
	boxrule = 0sp,
	borderline west = {2pt}{0pt}{myexercisefg},
	sharp corners,
	detach title,
	before upper = \tcbtitle\par\smallskip,
	coltitle = myexercisefg,
	fonttitle = \bfseries\sffamily,
	description font = \mdseries,
	separator sign none,
	segmentation style={solid, myexercisefg},
}
{th}

%================================
% EXAMPLE BOX
%================================

\newtcbtheorem[number within=section]{Example}{Example}
{%
	colback = myexamplebg
	,breakable
	,colframe = myexamplefr
	,coltitle = myexampleti
	,boxrule = 1pt
	,sharp corners
	,detach title
	,before upper=\tcbtitle\par\smallskip
	,fonttitle = \bfseries
	,description font = \mdseries
	,separator sign none
	,description delimiters parenthesis
}
{ex}

\newtcbtheorem[number within=chapter]{example}{Example}
{%
	colback = myexamplebg
	,breakable
	,colframe = myexamplefr
	,coltitle = myexampleti
	,boxrule = 1pt
	,sharp corners
	,detach title
	,before upper=\tcbtitle\par\smallskip
	,fonttitle = \bfseries
	,description font = \mdseries
	,separator sign none
	,description delimiters parenthesis
}
{ex}

%================================
% DEFINITION BOX
%================================

\newtcbtheorem[number within=section]{Definition}{Definition}{enhanced,
	before skip=2mm,after skip=2mm, colback=red!5,colframe=red!80!black,boxrule=0.5mm,
	attach boxed title to top left={xshift=1cm,yshift*=1mm-\tcboxedtitleheight}, varwidth boxed title*=-3cm,
	boxed title style={frame code={
					\path[fill=tcbcolback]
					([yshift=-1mm,xshift=-1mm]frame.north west)
					arc[start angle=0,end angle=180,radius=1mm]
					([yshift=-1mm,xshift=1mm]frame.north east)
					arc[start angle=180,end angle=0,radius=1mm];
					\path[left color=tcbcolback!60!black,right color=tcbcolback!60!black,
						middle color=tcbcolback!80!black]
					([xshift=-2mm]frame.north west) -- ([xshift=2mm]frame.north east)
					[rounded corners=1mm]-- ([xshift=1mm,yshift=-1mm]frame.north east)
					-- (frame.south east) -- (frame.south west)
					-- ([xshift=-1mm,yshift=-1mm]frame.north west)
					[sharp corners]-- cycle;
				},interior engine=empty,
		},
	fonttitle=\bfseries,
	title={#2},#1}{def}
\newtcbtheorem[number within=chapter]{definition}{Definition}{enhanced,
	before skip=2mm,after skip=2mm, colback=red!5,colframe=red!80!black,boxrule=0.5mm,
	attach boxed title to top left={xshift=1cm,yshift*=1mm-\tcboxedtitleheight}, varwidth boxed title*=-3cm,
	boxed title style={frame code={
					\path[fill=tcbcolback]
					([yshift=-1mm,xshift=-1mm]frame.north west)
					arc[start angle=0,end angle=180,radius=1mm]
					([yshift=-1mm,xshift=1mm]frame.north east)
					arc[start angle=180,end angle=0,radius=1mm];
					\path[left color=tcbcolback!60!black,right color=tcbcolback!60!black,
						middle color=tcbcolback!80!black]
					([xshift=-2mm]frame.north west) -- ([xshift=2mm]frame.north east)
					[rounded corners=1mm]-- ([xshift=1mm,yshift=-1mm]frame.north east)
					-- (frame.south east) -- (frame.south west)
					-- ([xshift=-1mm,yshift=-1mm]frame.north west)
					[sharp corners]-- cycle;
				},interior engine=empty,
		},
	fonttitle=\bfseries,
	title={#2},#1}{def}



%================================
% Solution BOX
%================================

\makeatletter
\newtcbtheorem{question}{Question}{enhanced,
	breakable,
	colback=white,
	colframe=myb!80!black,
	attach boxed title to top left={yshift*=-\tcboxedtitleheight},
	fonttitle=\bfseries,
	title={#2},
	boxed title size=title,
	boxed title style={%
			sharp corners,
			rounded corners=northwest,
			colback=tcbcolframe,
			boxrule=0pt,
		},
	underlay boxed title={%
			\path[fill=tcbcolframe] (title.south west)--(title.south east)
			to[out=0, in=180] ([xshift=5mm]title.east)--
			(title.center-|frame.east)
			[rounded corners=\kvtcb@arc] |-
			(frame.north) -| cycle;
		},
	#1
}{def}
\makeatother

%================================
% SOLUTION BOX
%================================

\makeatletter
\newtcolorbox{solution}{enhanced,
	breakable,
	colback=white,
	colframe=myg!80!black,
	attach boxed title to top left={yshift*=-\tcboxedtitleheight},
	title=Solution,
	boxed title size=title,
	boxed title style={%
			sharp corners,
			rounded corners=northwest,
			colback=tcbcolframe,
			boxrule=0pt,
		},
	underlay boxed title={%
			\path[fill=tcbcolframe] (title.south west)--(title.south east)
			to[out=0, in=180] ([xshift=5mm]title.east)--
			(title.center-|frame.east)
			[rounded corners=\kvtcb@arc] |-
			(frame.north) -| cycle;
		},
}
\makeatother

%================================
% Question BOX
%================================

\makeatletter
\newtcbtheorem{qstion}{Question}{enhanced,
	breakable,
	colback=white,
	colframe=mygr,
	attach boxed title to top left={yshift*=-\tcboxedtitleheight},
	fonttitle=\bfseries,
	title={#2},
	boxed title size=title,
	boxed title style={%
			sharp corners,
			rounded corners=northwest,
			colback=tcbcolframe,
			boxrule=0pt,
		},
	underlay boxed title={%
			\path[fill=tcbcolframe] (title.south west)--(title.south east)
			to[out=0, in=180] ([xshift=5mm]title.east)--
			(title.center-|frame.east)
			[rounded corners=\kvtcb@arc] |-
			(frame.north) -| cycle;
		},
	#1
}{def}
\makeatother

\newtcbtheorem[number within=chapter]{wconc}{Wrong Concept}{
	breakable,
	enhanced,
	colback=white,
	colframe=myr,
	arc=0pt,
	outer arc=0pt,
	fonttitle=\bfseries\sffamily\large,
	colbacktitle=myr,
	attach boxed title to top left={},
	boxed title style={
			enhanced,
			skin=enhancedfirst jigsaw,
			arc=3pt,
			bottom=0pt,
			interior style={fill=myr}
		},
	#1
}{def}



%================================
% NOTE BOX
%================================

\usetikzlibrary{arrows,calc,shadows.blur}
\tcbuselibrary{skins}
\newtcolorbox{note}[1][]{%
	enhanced jigsaw,
	colback=gray!20!white,%
	colframe=gray!80!black,
	size=small,
	boxrule=1pt,
	title=\textbf{Note:},
	halign title=flush center,
	coltitle=black,
	breakable,
	drop shadow=black!50!white,
	attach boxed title to top left={xshift=1cm,yshift=-\tcboxedtitleheight/2,yshifttext=-\tcboxedtitleheight/2},
	minipage boxed title=1.5cm,
	boxed title style={%
			colback=white,
			size=fbox,
			boxrule=1pt,
			boxsep=2pt,
			underlay={%
					\coordinate (dotA) at ($(interior.west) + (-0.5pt,0)$);
					\coordinate (dotB) at ($(interior.east) + (0.5pt,0)$);
					\begin{scope}
						\clip (interior.north west) rectangle ([xshift=3ex]interior.east);
						\filldraw [white, blur shadow={shadow opacity=60, shadow yshift=-.75ex}, rounded corners=2pt] (interior.north west) rectangle (interior.south east);
					\end{scope}
					\begin{scope}[gray!80!black]
						\fill (dotA) circle (2pt);
						\fill (dotB) circle (2pt);
					\end{scope}
				},
		},
	#1,
}

%%%%%%%%%%%%%%%%%%%%%%%%%%%%%%
% SELF MADE COMMANDS
%%%%%%%%%%%%%%%%%%%%%%%%%%%%%%


\newcommand{\thm}[2]{\begin{Theorem}{#1}{}#2\end{Theorem}}
\newcommand{\cor}[2]{\begin{Corollary}{#1}{}#2\end{Corollary}}
\newcommand{\mlenma}[2]{\begin{Lenma}{#1}{}#2\end{Lenma}}
\newcommand{\mprop}[2]{\begin{Prop}{#1}{}#2\end{Prop}}
\newcommand{\clm}[3]{\begin{claim}{#1}{#2}#3\end{claim}}
\newcommand{\wc}[2]{\begin{wconc}{#1}{}\setlength{\parindent}{1cm}#2\end{wconc}}
\newcommand{\thmcon}[1]{\begin{Theoremcon}{#1}\end{Theoremcon}}
\newcommand{\ex}[2]{\begin{Example}{#1}{}#2\end{Example}}
\newcommand{\dfn}[2]{\begin{Definition}[colbacktitle=red!75!black]{#1}{}#2\end{Definition}}
\newcommand{\dfnc}[2]{\begin{definition}[colbacktitle=red!75!black]{#1}{}#2\end{definition}}
\newcommand{\qs}[2]{\begin{question}{#1}{}#2\end{question}}
\newcommand{\pf}[2]{\begin{myproof}[#1]#2\end{myproof}}
\newcommand{\nt}[1]{\begin{note}#1\end{note}}

\newcommand*\circled[1]{\tikz[baseline=(char.base)]{
		\node[shape=circle,draw,inner sep=1pt] (char) {#1};}}
\newcommand\getcurrentref[1]{%
	\ifnumequal{\value{#1}}{0}
	{??}
	{\the\value{#1}}%
}
\newcommand{\getCurrentSectionNumber}{\getcurrentref{section}}
\newenvironment{myproof}[1][\proofname]{%
	\proof[\bfseries #1: ]%
}{\endproof}

\newcommand{\mclm}[2]{\begin{myclaim}[#1]#2\end{myclaim}}
\newenvironment{myclaim}[1][\claimname]{\proof[\bfseries #1: ]}{}

\newcounter{mylabelcounter}

\makeatletter
\newcommand{\setword}[2]{%
	\phantomsection
	#1\def\@currentlabel{\unexpanded{#1}}\label{#2}%
}
\makeatother




\tikzset{
	symbol/.style={
			draw=none,
			every to/.append style={
					edge node={node [sloped, allow upside down, auto=false]{$#1$}}}
		}
}


% deliminators
\DeclarePairedDelimiter{\abs}{\lvert}{\rvert}
\DeclarePairedDelimiter{\norm}{\lVert}{\rVert}

\DeclarePairedDelimiter{\ceil}{\lceil}{\rceil}
\DeclarePairedDelimiter{\floor}{\lfloor}{\rfloor}
\DeclarePairedDelimiter{\round}{\lfloor}{\rceil}

\newsavebox\diffdbox
\newcommand{\slantedromand}{{\mathpalette\makesl{d}}}
\newcommand{\makesl}[2]{%
\begingroup
\sbox{\diffdbox}{$\mathsurround=0pt#1\mathrm{#2}$}%
\pdfsave
\pdfsetmatrix{1 0 0.2 1}%
\rlap{\usebox{\diffdbox}}%
\pdfrestore
\hskip\wd\diffdbox
\endgroup
}
\newcommand{\dd}[1][]{\ensuremath{\mathop{}\!\ifstrempty{#1}{%
\slantedromand\@ifnextchar^{\hspace{0.2ex}}{\hspace{0.1ex}}}%
{\slantedromand\hspace{0.2ex}^{#1}}}}
\ProvideDocumentCommand\dv{o m g}{%
  \ensuremath{%
    \IfValueTF{#3}{%
      \IfNoValueTF{#1}{%
        \frac{\dd #2}{\dd #3}%
      }{%
        \frac{\dd^{#1} #2}{\dd #3^{#1}}%
      }%
    }{%
      \IfNoValueTF{#1}{%
        \frac{\dd}{\dd #2}%
      }{%
        \frac{\dd^{#1}}{\dd #2^{#1}}%
      }%
    }%
  }%
}
\providecommand*{\pdv}[3][]{\frac{\partial^{#1}#2}{\partial#3^{#1}}}
%  - others
\DeclareMathOperator{\Lap}{\mathcal{L}}
\DeclareMathOperator{\Var}{Var} % varience
\DeclareMathOperator{\Cov}{Cov} % covarience
\DeclareMathOperator{\E}{E} % expected

% Since the amsthm package isn't loaded

% I prefer the slanted \leq
\let\oldleq\leq % save them in case they're every wanted
\let\oldgeq\geq
\renewcommand{\leq}{\leqslant}
\renewcommand{\geq}{\geqslant}

% % redefine matrix env to allow for alignment, use r as default
% \renewcommand*\env@matrix[1][r]{\hskip -\arraycolsep
%     \let\@ifnextchar\new@ifnextchar
%     \array{*\c@MaxMatrixCols #1}}


%\usepackage{framed}
%\usepackage{titletoc}
%\usepackage{etoolbox}
%\usepackage{lmodern}


%\patchcmd{\tableofcontents}{\contentsname}{\sffamily\contentsname}{}{}

%\renewenvironment{leftbar}
%{\def\FrameCommand{\hspace{6em}%
%		{\color{myyellow}\vrule width 2pt depth 6pt}\hspace{1em}}%
%	\MakeFramed{\parshape 1 0cm \dimexpr\textwidth-6em\relax\FrameRestore}\vskip2pt%
%}
%{\endMakeFramed}

%\titlecontents{chapter}
%[0em]{\vspace*{2\baselineskip}}
%{\parbox{4.5em}{%
%		\hfill\Huge\sffamily\bfseries\color{myred}\thecontentspage}%
%	\vspace*{-2.3\baselineskip}\leftbar\textsc{\small\chaptername~\thecontentslabel}\\\sffamily}
%{}{\endleftbar}
%\titlecontents{section}
%[8.4em]
%{\sffamily\contentslabel{3em}}{}{}
%{\hspace{0.5em}\nobreak\itshape\color{myred}\contentspage}
%\titlecontents{subsection}
%[8.4em]
%{\sffamily\contentslabel{3em}}{}{}  
%{\hspace{0.5em}\nobreak\itshape\color{myred}\contentspage}



%%%%%%%%%%%%%%%%%%%%%%%%%%%%%%%%%%%%%%%%%%%
% TABLE OF CONTENTS
%%%%%%%%%%%%%%%%%%%%%%%%%%%%%%%%%%%%%%%%%%%

\usepackage{tikz}
\definecolor{doc}{RGB}{0,60,110}
\usepackage{titletoc}
\contentsmargin{0cm}
\titlecontents{chapter}[3.7pc]
{\addvspace{30pt}%
	\begin{tikzpicture}[remember picture, overlay]%
		\draw[fill=doc!60,draw=doc!60] (-7,-.1) rectangle (-0.9,.5);%
		\pgftext[left,x=-3.5cm,y=0.2cm]{\color{white}\Large\sc\bfseries Chapter\ \thecontentslabel};%
	\end{tikzpicture}\color{doc!60}\large\sc\bfseries}%
{}
{}
{\;\titlerule\;\large\sc\bfseries Page \thecontentspage
	\begin{tikzpicture}[remember picture, overlay]
		\draw[fill=doc!60,draw=doc!60] (2pt,0) rectangle (4,0.1pt);
	\end{tikzpicture}}%
\titlecontents{section}[3.7pc]
{\addvspace{2pt}}
{\contentslabel[\thecontentslabel]{2pc}}
{}
{\hfill\small \thecontentspage}
[]
\titlecontents*{subsection}[3.7pc]
{\addvspace{-1pt}\small}
{}
{}
{\ --- \small\thecontentspage}
[ \textbullet\ ][]

\makeatletter
\renewcommand{\tableofcontents}{%
	\chapter*{%
	  \vspace*{-20\p@}%
	  \begin{tikzpicture}[remember picture, overlay]%
		  \pgftext[right,x=15cm,y=0.2cm]{\color{doc!60}\Huge\sc\bfseries \contentsname};%
		  \draw[fill=doc!60,draw=doc!60] (13,-.75) rectangle (20,1);%
		  \clip (13,-.75) rectangle (20,1);
		  \pgftext[right,x=15cm,y=0.2cm]{\color{white}\Huge\sc\bfseries \contentsname};%
	  \end{tikzpicture}}%
	\@starttoc{toc}}
\makeatother


% Comandi di Bastiality %
\newcommand{\innerproduct}[2]{\langle #1, #2 \rangle}
\newcommand{\generators}[1]{\langle #1 \rangle}
\newcommand{\powerset}[0]{\mathcal{P}}
\newcommand{\indip}{\perp \!\!\! \perp}

% Comandi di GioLaPalma %
\newcommand{\der}[0]{\implies^*}
\newcommand{\red}[1]{\textcolor{red}{#1}}
\newcommand{\esempio}[1]{\ex{}{#1}}
\newcommand{\teorema}[1]{\thm{}{#1}}
\newcommand{\dimostrazione}[1]{\pf{Dimostrazione}{#1}}
\setlength{\parindent}{0pt}
\newcommand{\osservazione}[1]{\nt{#1}}
\newcommand{\inside}[1]{
  \langle #1 \rangle
}
      

%From M275 "Topology" at SJSU
\newcommand{\id}{\mathrm{id}}
\newcommand{\taking}[1]{\xrightarrow{#1}}
\newcommand{\inv}{^{-1}}

%From M170 "Introduction to Graph Theory" at SJSU
\DeclareMathOperator{\diam}{diam}
\DeclareMathOperator{\ord}{ord}
\newcommand{\defeq}{\overset{\mathrm{def}}{=}}

%From the USAMO .tex files
\newcommand{\ts}{\textsuperscript}
\newcommand{\dg}{^\circ}
\newcommand{\ii}{\item}

% % From Math 55 and Math 145 at Harvard
% \newenvironment{subproof}[1][Proof]{%
% \begin{proof}[#1] \renewcommand{\qedsymbol}{$\blacksquare$}}%
% {\end{proof}}

\newcommand{\liff}{\leftrightarrow}
\newcommand{\lthen}{\rightarrow}
\newcommand{\opname}{\operatorname}
\newcommand{\surjto}{\twoheadrightarrow}
\newcommand{\injto}{\hookrightarrow}
\newcommand{\On}{\mathrm{On}} % ordinals
\DeclareMathOperator{\img}{im} % Image
\DeclareMathOperator{\Img}{Im} % Image
\DeclareMathOperator{\coker}{coker} % Cokernel
\DeclareMathOperator{\Coker}{Coker} % Cokernel
\DeclareMathOperator{\Ker}{Ker} % Kernel
\DeclareMathOperator{\rank}{rank}
\DeclareMathOperator{\Spec}{Spec} % spectrum
\DeclareMathOperator{\Tr}{Tr} % trace
\DeclareMathOperator{\pr}{pr} % projection
\DeclareMathOperator{\ext}{ext} % extension
\DeclareMathOperator{\pred}{pred} % predecessor
\DeclareMathOperator{\dom}{dom} % domain
\DeclareMathOperator{\ran}{ran} % range
\DeclareMathOperator{\Hom}{Hom} % homomorphism
\DeclareMathOperator{\Mor}{Mor} % morphisms
\DeclareMathOperator{\End}{End} % endomorphism

\newcommand{\eps}{\epsilon}
\newcommand{\veps}{\varepsilon}
\newcommand{\ol}{\overline}
\newcommand{\ul}{\underline}
\newcommand{\wt}{\widetilde}
\newcommand{\wh}{\widehat}
\newcommand{\vocab}[1]{\textbf{\color{blue} #1}}
\providecommand{\half}{\frac{1}{2}}
\newcommand{\dang}{\measuredangle} %% Directed angle
\newcommand{\ray}[1]{\overrightarrow{#1}}
\newcommand{\seg}[1]{\overline{#1}}
\newcommand{\arc}[1]{\wideparen{#1}}
\DeclareMathOperator{\cis}{cis}
\DeclareMathOperator*{\lcm}{lcm}
\DeclareMathOperator*{\argmin}{arg min}
\DeclareMathOperator*{\argmax}{arg max}
\newcommand{\cycsum}{\sum_{\mathrm{cyc}}}
\newcommand{\symsum}{\sum_{\mathrm{sym}}}
\newcommand{\cycprod}{\prod_{\mathrm{cyc}}}
\newcommand{\symprod}{\prod_{\mathrm{sym}}}
\newcommand{\Qed}{\begin{flushright}\qed\end{flushright}}
\newcommand{\parinn}{\setlength{\parindent}{1cm}}
\newcommand{\parinf}{\setlength{\parindent}{0cm}}
% \newcommand{\norm}{\|\cdot\|}
\newcommand{\inorm}{\norm_{\infty}}
\newcommand{\opensets}{\{V_{\alpha}\}_{\alpha\in I}}
\newcommand{\oset}{V_{\alpha}}
\newcommand{\opset}[1]{V_{\alpha_{#1}}}
\newcommand{\lub}{\text{lub}}
\newcommand{\del}[2]{\frac{\partial #1}{\partial #2}}
\newcommand{\Del}[3]{\frac{\partial^{#1} #2}{\partial^{#1} #3}}
\newcommand{\deld}[2]{\dfrac{\partial #1}{\partial #2}}
\newcommand{\Deld}[3]{\dfrac{\partial^{#1} #2}{\partial^{#1} #3}}
\newcommand{\lm}{\lambda}
\newcommand{\uin}{\mathbin{\rotatebox[origin=c]{90}{$\in$}}}
\newcommand{\usubset}{\mathbin{\rotatebox[origin=c]{90}{$\subset$}}}
\newcommand{\lt}{\left}
\newcommand{\rt}{\right}
\newcommand{\bs}[1]{\boldsymbol{#1}}
\newcommand{\exs}{\exists}
\newcommand{\st}{\strut}
\newcommand{\dps}[1]{\displaystyle{#1}}

\newcommand{\sol}{\setlength{\parindent}{0cm}\textbf{\textit{Solution:}}\setlength{\parindent}{1cm} }
\newcommand{\solve}[1]{\setlength{\parindent}{0cm}\textbf{\textit{Solution: }}\setlength{\parindent}{1cm}#1 \Qed}

% Things Lie
\newcommand{\kb}{\mathfrak b}
\newcommand{\kg}{\mathfrak g}
\newcommand{\kh}{\mathfrak h}
\newcommand{\kn}{\mathfrak n}
\newcommand{\ku}{\mathfrak u}
\newcommand{\kz}{\mathfrak z}
\DeclareMathOperator{\Ext}{Ext} % Ext functor
\DeclareMathOperator{\Tor}{Tor} % Tor functor
\newcommand{\gl}{\opname{\mathfrak{gl}}} % frak gl group
\renewcommand{\sl}{\opname{\mathfrak{sl}}} % frak sl group chktex 6

% More script letters etc.
\newcommand{\SA}{\mathcal A}
\newcommand{\SB}{\mathcal B}
\newcommand{\SC}{\mathcal C}
\newcommand{\SF}{\mathcal F}
\newcommand{\SG}{\mathcal G}
\newcommand{\SH}{\mathcal H}
\newcommand{\OO}{\mathcal O}

\newcommand{\SCA}{\mathscr A}
\newcommand{\SCB}{\mathscr B}
\newcommand{\SCC}{\mathscr C}
\newcommand{\SCD}{\mathscr D}
\newcommand{\SCE}{\mathscr E}
\newcommand{\SCF}{\mathscr F}
\newcommand{\SCG}{\mathscr G}
\newcommand{\SCH}{\mathscr H}

% Mathfrak primes
\newcommand{\km}{\mathfrak m}
\newcommand{\kp}{\mathfrak p}
\newcommand{\kq}{\mathfrak q}

% number sets
\newcommand{\RR}[1][]{\ensuremath{\ifstrempty{#1}{\mathbb{R}}{\mathbb{R}^{#1}}}}
\newcommand{\NN}[1][]{\ensuremath{\ifstrempty{#1}{\mathbb{N}}{\mathbb{N}^{#1}}}}
\newcommand{\ZZ}[1][]{\ensuremath{\ifstrempty{#1}{\mathbb{Z}}{\mathbb{Z}^{#1}}}}
\newcommand{\QQ}[1][]{\ensuremath{\ifstrempty{#1}{\mathbb{Q}}{\mathbb{Q}^{#1}}}}
\newcommand{\CC}[1][]{\ensuremath{\ifstrempty{#1}{\mathbb{C}}{\mathbb{C}^{#1}}}}
\newcommand{\PP}[1][]{\ensuremath{\ifstrempty{#1}{\mathbb{P}}{\mathbb{P}^{#1}}}}
\newcommand{\HH}[1][]{\ensuremath{\ifstrempty{#1}{\mathbb{H}}{\mathbb{H}^{#1}}}}
\newcommand{\FF}[1][]{\ensuremath{\ifstrempty{#1}{\mathbb{F}}{\mathbb{F}^{#1}}}}
% expected value
\newcommand{\EE}{\ensuremath{\mathbb{E}}}
\newcommand{\charin}{\text{ char }}
\DeclareMathOperator{\sign}{sign}
\DeclareMathOperator{\Aut}{Aut}
\DeclareMathOperator{\Inn}{Inn}
\DeclareMathOperator{\Syl}{Syl}
\DeclareMathOperator{\Gal}{Gal}
\DeclareMathOperator{\GL}{GL} % General linear group
\DeclareMathOperator{\SL}{SL} % Special linear group

%---------------------------------------
% BlackBoard Math Fonts :-
%---------------------------------------

%Captital Letters
\newcommand{\bbA}{\mathbb{A}}	\newcommand{\bbB}{\mathbb{B}}
\newcommand{\bbC}{\mathbb{C}}	\newcommand{\bbD}{\mathbb{D}}
\newcommand{\bbE}{\mathbb{E}}	\newcommand{\bbF}{\mathbb{F}}
\newcommand{\bbG}{\mathbb{G}}	\newcommand{\bbH}{\mathbb{H}}
\newcommand{\bbI}{\mathbb{I}}	\newcommand{\bbJ}{\mathbb{J}}
\newcommand{\bbK}{\mathbb{K}}	\newcommand{\bbL}{\mathbb{L}}
\newcommand{\bbM}{\mathbb{M}}	\newcommand{\bbN}{\mathbb{N}}
\newcommand{\bbO}{\mathbb{O}}	\newcommand{\bbP}{\mathbb{P}}
\newcommand{\bbQ}{\mathbb{Q}}	\newcommand{\bbR}{\mathbb{R}}
\newcommand{\bbS}{\mathbb{S}}	\newcommand{\bbT}{\mathbb{T}}
\newcommand{\bbU}{\mathbb{U}}	\newcommand{\bbV}{\mathbb{V}}
\newcommand{\bbW}{\mathbb{W}}	\newcommand{\bbX}{\mathbb{X}}
\newcommand{\bbY}{\mathbb{Y}}	\newcommand{\bbZ}{\mathbb{Z}}

%---------------------------------------
% MathCal Fonts :-
%---------------------------------------

%Captital Letters
\newcommand{\mcA}{\mathcal{A}}	\newcommand{\mcB}{\mathcal{B}}
\newcommand{\mcC}{\mathcal{C}}	\newcommand{\mcD}{\mathcal{D}}
\newcommand{\mcE}{\mathcal{E}}	\newcommand{\mcF}{\mathcal{F}}
\newcommand{\mcG}{\mathcal{G}}	\newcommand{\mcH}{\mathcal{H}}
\newcommand{\mcI}{\mathcal{I}}	\newcommand{\mcJ}{\mathcal{J}}
\newcommand{\mcK}{\mathcal{K}}	\newcommand{\mcL}{\mathcal{L}}
\newcommand{\mcM}{\mathcal{M}}	\newcommand{\mcN}{\mathcal{N}}
\newcommand{\mcO}{\mathcal{O}}	\newcommand{\mcP}{\mathcal{P}}
\newcommand{\mcQ}{\mathcal{Q}}	\newcommand{\mcR}{\mathcal{R}}
\newcommand{\mcS}{\mathcal{S}}	\newcommand{\mcT}{\mathcal{T}}
\newcommand{\mcU}{\mathcal{U}}	\newcommand{\mcV}{\mathcal{V}}
\newcommand{\mcW}{\mathcal{W}}	\newcommand{\mcX}{\mathcal{X}}
\newcommand{\mcY}{\mathcal{Y}}	\newcommand{\mcZ}{\mathcal{Z}}


%---------------------------------------
% Bold Math Fonts :-
%---------------------------------------

%Captital Letters
\newcommand{\bmA}{\boldsymbol{A}}	\newcommand{\bmB}{\boldsymbol{B}}
\newcommand{\bmC}{\boldsymbol{C}}	\newcommand{\bmD}{\boldsymbol{D}}
\newcommand{\bmE}{\boldsymbol{E}}	\newcommand{\bmF}{\boldsymbol{F}}
\newcommand{\bmG}{\boldsymbol{G}}	\newcommand{\bmH}{\boldsymbol{H}}
\newcommand{\bmI}{\boldsymbol{I}}	\newcommand{\bmJ}{\boldsymbol{J}}
\newcommand{\bmK}{\boldsymbol{K}}	\newcommand{\bmL}{\boldsymbol{L}}
\newcommand{\bmM}{\boldsymbol{M}}	\newcommand{\bmN}{\boldsymbol{N}}
\newcommand{\bmO}{\boldsymbol{O}}	\newcommand{\bmP}{\boldsymbol{P}}
\newcommand{\bmQ}{\boldsymbol{Q}}	\newcommand{\bmR}{\boldsymbol{R}}
\newcommand{\bmS}{\boldsymbol{S}}	\newcommand{\bmT}{\boldsymbol{T}}
\newcommand{\bmU}{\boldsymbol{U}}	\newcommand{\bmV}{\boldsymbol{V}}
\newcommand{\bmW}{\boldsymbol{W}}	\newcommand{\bmX}{\boldsymbol{X}}
\newcommand{\bmY}{\boldsymbol{Y}}	\newcommand{\bmZ}{\boldsymbol{Z}}
%Small Letters
\newcommand{\bma}{\boldsymbol{a}}	\newcommand{\bmb}{\boldsymbol{b}}
\newcommand{\bmc}{\boldsymbol{c}}	\newcommand{\bmd}{\boldsymbol{d}}
\newcommand{\bme}{\boldsymbol{e}}	\newcommand{\bmf}{\boldsymbol{f}}
\newcommand{\bmg}{\boldsymbol{g}}	\newcommand{\bmh}{\boldsymbol{h}}
\newcommand{\bmi}{\boldsymbol{i}}	\newcommand{\bmj}{\boldsymbol{j}}
\newcommand{\bmk}{\boldsymbol{k}}	\newcommand{\bml}{\boldsymbol{l}}
\newcommand{\bmm}{\boldsymbol{m}}	\newcommand{\bmn}{\boldsymbol{n}}
\newcommand{\bmo}{\boldsymbol{o}}	\newcommand{\bmp}{\boldsymbol{p}}
\newcommand{\bmq}{\boldsymbol{q}}	\newcommand{\bmr}{\boldsymbol{r}}
\newcommand{\bms}{\boldsymbol{s}}	\newcommand{\bmt}{\boldsymbol{t}}
\newcommand{\bmu}{\boldsymbol{u}}	\newcommand{\bmv}{\boldsymbol{v}}
\newcommand{\bmw}{\boldsymbol{w}}	\newcommand{\bmx}{\boldsymbol{x}}
\newcommand{\bmy}{\boldsymbol{y}}	\newcommand{\bmz}{\boldsymbol{z}}

%---------------------------------------
% Scr Math Fonts :-
%---------------------------------------

\newcommand{\sA}{{\mathscr{A}}}   \newcommand{\sB}{{\mathscr{B}}}
\newcommand{\sC}{{\mathscr{C}}}   \newcommand{\sD}{{\mathscr{D}}}
\newcommand{\sE}{{\mathscr{E}}}   \newcommand{\sF}{{\mathscr{F}}}
\newcommand{\sG}{{\mathscr{G}}}   \newcommand{\sH}{{\mathscr{H}}}
\newcommand{\sI}{{\mathscr{I}}}   \newcommand{\sJ}{{\mathscr{J}}}
\newcommand{\sK}{{\mathscr{K}}}   \newcommand{\sL}{{\mathscr{L}}}
\newcommand{\sM}{{\mathscr{M}}}   \newcommand{\sN}{{\mathscr{N}}}
\newcommand{\sO}{{\mathscr{O}}}   \newcommand{\sP}{{\mathscr{P}}}
\newcommand{\sQ}{{\mathscr{Q}}}   \newcommand{\sR}{{\mathscr{R}}}
\newcommand{\sS}{{\mathscr{S}}}   \newcommand{\sT}{{\mathscr{T}}}
\newcommand{\sU}{{\mathscr{U}}}   \newcommand{\sV}{{\mathscr{V}}}
\newcommand{\sW}{{\mathscr{W}}}   \newcommand{\sX}{{\mathscr{X}}}
\newcommand{\sY}{{\mathscr{Y}}}   \newcommand{\sZ}{{\mathscr{Z}}}


%---------------------------------------
% Math Fraktur Font
%---------------------------------------

%Captital Letters
\newcommand{\mfA}{\mathfrak{A}}	\newcommand{\mfB}{\mathfrak{B}}
\newcommand{\mfC}{\mathfrak{C}}	\newcommand{\mfD}{\mathfrak{D}}
\newcommand{\mfE}{\mathfrak{E}}	\newcommand{\mfF}{\mathfrak{F}}
\newcommand{\mfG}{\mathfrak{G}}	\newcommand{\mfH}{\mathfrak{H}}
\newcommand{\mfI}{\mathfrak{I}}	\newcommand{\mfJ}{\mathfrak{J}}
\newcommand{\mfK}{\mathfrak{K}}	\newcommand{\mfL}{\mathfrak{L}}
\newcommand{\mfM}{\mathfrak{M}}	\newcommand{\mfN}{\mathfrak{N}}
\newcommand{\mfO}{\mathfrak{O}}	\newcommand{\mfP}{\mathfrak{P}}
\newcommand{\mfQ}{\mathfrak{Q}}	\newcommand{\mfR}{\mathfrak{R}}
\newcommand{\mfS}{\mathfrak{S}}	\newcommand{\mfT}{\mathfrak{T}}
\newcommand{\mfU}{\mathfrak{U}}	\newcommand{\mfV}{\mathfrak{V}}
\newcommand{\mfW}{\mathfrak{W}}	\newcommand{\mfX}{\mathfrak{X}}
\newcommand{\mfY}{\mathfrak{Y}}	\newcommand{\mfZ}{\mathfrak{Z}}
%Small Letters
\newcommand{\mfa}{\mathfrak{a}}	\newcommand{\mfb}{\mathfrak{b}}
\newcommand{\mfc}{\mathfrak{c}}	\newcommand{\mfd}{\mathfrak{d}}
\newcommand{\mfe}{\mathfrak{e}}	\newcommand{\mff}{\mathfrak{f}}
\newcommand{\mfg}{\mathfrak{g}}	\newcommand{\mfh}{\mathfrak{h}}
\newcommand{\mfi}{\mathfrak{i}}	\newcommand{\mfj}{\mathfrak{j}}
\newcommand{\mfk}{\mathfrak{k}}	\newcommand{\mfl}{\mathfrak{l}}
\newcommand{\mfm}{\mathfrak{m}}	\newcommand{\mfn}{\mathfrak{n}}
\newcommand{\mfo}{\mathfrak{o}}	\newcommand{\mfp}{\mathfrak{p}}
\newcommand{\mfq}{\mathfrak{q}}	\newcommand{\mfr}{\mathfrak{r}}
\newcommand{\mfs}{\mathfrak{s}}	\newcommand{\mft}{\mathfrak{t}}
\newcommand{\mfu}{\mathfrak{u}}	\newcommand{\mfv}{\mathfrak{v}}
\newcommand{\mfw}{\mathfrak{w}}	\newcommand{\mfx}{\mathfrak{x}}
\newcommand{\mfy}{\mathfrak{y}}	\newcommand{\mfz}{\mathfrak{z}}


\usepackage[T1]{fontenc}
\usepackage[utf8]{inputenc}
\usepackage{xcolor}
\usepackage{listings}
\usepackage{array}
\usepackage{longtable}
\PassOptionsToPackage{hyphens}{url}
\usepackage{url}
\Urlmuskip=0mu plus 1mu

%==============================================================
\definecolor{codebg}{HTML}{F7F7FB}
\definecolor{codekw}{HTML}{2D6FB1}
\definecolor{codecm}{HTML}{388C46}
\definecolor{codestr}{HTML}{C74440}
\definecolor{coderule}{HTML}{CCCCCC}

\lstset{
    basicstyle=\ttfamily\footnotesize,
    keywordstyle=\color{codekw}\bfseries,
    commentstyle=\color{codecm}\itshape,
    stringstyle=\color{codestr},
    numberstyle=\tiny\color{gray},
    showstringspaces=false,
    breaklines=true,
    breakatwhitespace=false,
    frame=single,
    rulecolor=\color{coderule},
    backgroundcolor=\color{codebg},
    tabsize=2,
    columns=fullflexible,
    keepspaces=true,
    captionpos=b,
    xleftmargin=1.2em,
    aboveskip=0.9em,
    belowskip=0.9em,
}

\newcommand{\code}[1]{\texttt{\small #1}}
\newcommand{\fonteslide}[1]{{\footnotesize\itshape\color{gray} Fonte: slide \emph{ExtraLectureOnVF}, pag.~#1.\par}}
\newcommand{\fontewiki}[1][Virtual function]{{\footnotesize\itshape\color{gray} Fonte: Wikipedia, \emph{#1}.\par}}
\newcommand{\fontelibro}[1]{{\footnotesize\itshape\color{gray} Fonte: \emph{Thinking in C++}, Vol.~2, #1.\par}}
\newcommand{\fonteprove}[1]{{\footnotesize\itshape\color{gray} Verificato in \code{prove/#1} (compilato ed eseguito).\par}}

\graphicspath{{./img/}}

%==============================================================
\begin{document}

\begin{titlepage}
    \centering
    \vspace*{2.3cm}
    {\Huge\sffamily\bfseries Virtual Functions\par}
    \vspace{0.3cm}
    {\LARGE\sffamily e Ereditarietà Virtuale\par}
    \vspace{0.5cm}
    {\LARGE\sffamily Teoria completa\par}
    \vspace{1.2cm}
    {\Large Appunti di Ingegneria del Software\par}
    \vspace{0.4cm}
    {\large Basato sulle slide del corso (\emph{Virtual Functions and Late
    Bindings}), su Wikipedia (\emph{Virtual function}, \emph{Virtual
    inheritance}) e su \emph{Thinking in C++} (Bruce Eckel)\par}
    \vspace{1.6cm}
    {\large\itshape Late binding, overriding, VFT/VFTP, dominanza,\\
    diamond problem, basi virtuali\par}
    \vspace{0.6cm}
    {\footnotesize\itshape Ogni esempio di codice di questo documento è
    stato compilato ed eseguito davvero (cartella \code{prove/}): dove il
    comportamento reale si è rivelato diverso da quello ``da manuale'',
    è segnalato esplicitamente.\par}
    \vfill
    {\large Anno 3\par}
\end{titlepage}

\tableofcontents
\newpage

%==============================================================
\chapter{Perché servono le funzioni virtuali}\label{ch:intro}
%==============================================================

\section{Il problema: comportamento deciso il più tardi possibile}

Una delle idee centrali della programmazione a oggetti è che il
comportamento di un oggetto venga determinato \textbf{al momento più
tardo possibile}, se possibile a runtime. Nei linguaggi imperativi
convenzionali, questo era di solito possibile solo tramite istruzioni
\code{switch} o \code{if}. Le cosiddette \textbf{funzioni virtuali} sono
una soluzione elegante e flessibile a questo problema: implementano un
meccanismo chiamato \textbf{late binding}.

\fonteslide{3}

\dfn{Funzione virtuale}{
    Un metodo (non statico) di una classe base che può essere
    ``ridefinito'' nelle classi derivate. Questa ridefinizione si chiama
    \textbf{overriding}.
}

\fonteslide{5}

\section{Overriding: le condizioni}

L'overriding avviene solo se:
\begin{itemize}
    \item il metodo sovrascritto (nella base) è \textbf{virtual};
    \item il metodo che sovrascrive (nella derivata) ha la \textbf{stessa
          signature} e lo \textbf{stesso return type} del metodo
          sovrascritto.
\end{itemize}
L'overriding può essere impedito in certi linguaggi di programmazione.

\fonteslide{6}

\dfn{Signature vs. prototype}{
    La \textbf{signature} di una funzione include: nome, numero di
    parametri formali, tipi dei parametri formali. Il \textbf{return type
    non è incluso} nella signature. Signature + return type insieme si
    chiamano il \textbf{prototype} della funzione.
}

\fonteslide{7}

\mprop{Conseguenze pratiche delle regole di overriding}{
    \begin{itemize}
        \item i metodi \textbf{static non possono essere overridden};
        \item stessa signature ma return type diverso $\to$ \textbf{errore
              di compilazione} (il C++ ammette solo lievi variazioni
              covarianti del return type);
        \item in \textbf{Java}, stesso nome ma signature diversa $\to$
              \textbf{l'overriding non avviene} (overloading, non
              override);
        \item in \textbf{C++}, stesso nome ma signature diversa $\to$ il
              metodo con quel nome nella base diventa \textbf{invisibile}
              (hiding) nella derivata. L'overloading può avvenire o meno a
              seconda del linguaggio.
    \end{itemize}
}

\fonteslide{8-9}
\fonteprove{03\_name\_hiding\_overload.cpp}

\section{Le tre possibilità del dispatch}

Quando una funzione virtuale \code{f} viene invocata su un oggetto
\code{o} di classe \code{D} derivata da \code{B}, ci sono tre
possibilità: (1) \code{f} è definita solo in \code{D} $\to$ si chiama
\code{D.f}; (2) \code{f} è definita solo in \code{B} $\to$ si chiama
\code{B.f} (ereditata); (3) \code{f} è overridden in \code{D} $\to$ si
chiama la versione overridden \code{D.f}.

\mprop{La regola fondamentale}{
    Si usa sempre la \textbf{classe reale} dell'oggetto, \underline{non}
    la classe della variabile/riferimento/puntatore usato per accedervi.
}

\fonteslide{11}
\fonteprove{01\_dispatch\_realtype.cpp}

Il vero vantaggio delle funzioni virtuali emerge quando si manipolano
oggetti tramite \textbf{riferimenti o puntatori alla loro classe base}:
solo in questo caso il comportamento ``virtuale'' diventa visibile
all'utente.

\fonteslide{12}

%==============================================================
\chapter{Virtual functions in Java e in C++}\label{ch:javacpp}
%==============================================================

\section{Java: tutto è virtual di default}

In Java \textbf{tutti} i metodi non statici sono virtual: l'overriding è
sempre un'opzione. Per impedire l'overriding (e l'hiding) di un metodo si
usa la keyword \code{final}; dichiarare \code{final} un'intera classe ne
impedisce l'estensione (non può essere usata come base per ulteriore
derivazione).

\fonteslide{16}

\begin{Example}{L'esempio Animal/Dog/Cat delle slide}{}
\begin{lstlisting}[language=Java]
public class Animal {
    public void voice() { System.out.println("???"); }
}
public class Dog extends Animal {
    public void voice() { System.out.println("bau"); }
}
public class Cat extends Animal {
    public void voice() { System.out.println("miao"); }
}
public class VirtualF {
    public static void main(String argv[]) {
        Animal a = new Dog();
        Animal b = new Cat();
        a.voice();   // both a and b are references to "Animals"
        b.voice();
    }
}
\end{lstlisting}
Output: \code{bau} (\code{a} fa riferimento a un \code{Dog}), \code{miao}
(\code{b} a un \code{Cat}) --- il dispatch usa sempre la classe reale.
\end{Example}

\fonteslide{25}
\fonteprove{10\_VirtualFunctionsInJava.java}

\section{C++: virtual va dichiarato esplicitamente}

In C++ un metodo deve essere dichiarato \code{virtual} esplicitamente.
Un metodo \textbf{non} dichiarato \code{virtual} \textbf{non può essere
overridden}: può solo essere \textbf{nascosto} da una nuova definizione.
Non esiste, in C++, una keyword equivalente a \code{final} di Java: non è
possibile impedire l'hiding di un metodo. La derivazione di classe è
invece sempre possibile.

\fonteslide{18}

\begin{lstlisting}[language=C++]
class B {
public:
    virtual int compute(int x) { return theInt + x; }
};
class D : public B {
public:
    int compute(int x) { return x * theInt; }   // 'virtual' puo' essere omesso
};
\end{lstlisting}

\fonteslide{20}
\fonteprove{02\_hiding\_non\_virtual.cpp}

\wc{Se scrivo una funzione con lo stesso nome in una derivata, sto sempre facendo l'override}{
    No: se la funzione della base \textbf{non} è \code{virtual}, quella
    che scrivi nella derivata \textbf{nasconde} la versione della base,
    non la sovrascrive. Late binding non viene mai eseguito per quel
    metodo (esempio 02): chiamandolo tramite un puntatore/riferimento alla
    base, si esegue sempre la versione della base, indipendentemente
    dalla classe reale dell'oggetto.
}

%==============================================================
\chapter{Come si implementa il late binding}\label{ch:vtable}
%==============================================================

\section{Early binding vs late binding}

\dfn{Early binding (compile-time binding)}{
    Nei linguaggi imperativi non-OO (C, Pascal...), quando il compilatore
    incontra una chiamata a funzione scrive codice macchina che esegue
    una chiamata a un \textbf{indirizzo ben definito} --- l'indirizzo di
    memoria dove inizia il codice macchina della funzione. Il compilatore
    ``lega'' il punto di chiamata al codice giusto già a \textbf{compile
    time}. È il comportamento standard per funzioni e metodi C++ non
    virtual.
}

\fonteslide{29-30}

\dfn{Late binding (run-time binding)}{
    In presenza di funzioni virtuali, il binding \textbf{non può} essere
    fatto a compile time: il compilatore sa quale funzione è virtual, ma
    non può sempre dedurre quale codice eseguire guardando solo la
    chiamata. Scrive invece codice macchina che deduce l'indirizzo
    corretto della chiamata \textbf{a runtime}.
}

\fonteslide{31}

\section{VFT e VFTP: l'implementazione}

Ogni classe con funzioni virtuali ha una tabella associata, chiamata
\textbf{Virtual Function Table (VFT)}, contenente i puntatori al codice
di ciascuna funzione virtuale della classe. Ogni oggetto di una classe
con funzioni virtuali porta un puntatore, \textbf{nascosto all'utente},
alla VFT della propria classe: il \textbf{Virtual Function Table Pointer
(VFTP)}.

\fonteslide{33}

\nt{
    In questo contesto ``pointer'' significa sia un puntatore C++ sia un
    riferimento Java: il meccanismo è concettualmente lo stesso nei due
    linguaggi, anche se Java lo nasconde completamente.
}

\mprop{I cinque passi del dispatch a runtime}{
    Per completare una chiamata late-bound:
    \begin{enumerate}
        \item dal puntatore/riferimento si raggiunge l'\textbf{oggetto};
        \item si cerca la \textbf{VFTP} nell'oggetto;
        \item tramite la VFTP si trova la \textbf{VFT della classe reale}
              dell'oggetto;
        \item si cerca nella VFT l'\textbf{entry giusta} (es. VFP2 per la
              funzione VF2);
        \item tramite il puntatore in quella entry si raggiunge ed
              \textbf{esegue} il codice corretto.
    \end{enumerate}
}

\fonteslide{32,34-35}

Questa è esattamente la traduzione in C che si può scrivere a mano
(struct con un puntatore a una tabella di puntatori a funzione): vedi il
capitolo sui template per il parallelo con la genericità in C.

%==============================================================
\chapter{Virtual vs. non-virtual: l'efficienza}\label{ch:efficienza}
%==============================================================

Se la ``virtualità'' è così comoda, perché non usare sempre metodi
virtual? Java prende questo approccio, vincolando l'utente a usare solo
metodi virtual. Il C++ dà più libertà, permettendo anche metodi non
virtual --- ma perché?

\fonteslide{36}

\mprop{La risposta: efficienza}{
    Il late binding è un meccanismo relativamente costoso in termini di
    tempo di esecuzione (in misura diversa secondo l'implementazione
    usata). Se un metodo viene chiamato molto frequentemente ed è
    relativamente semplice, l'overhead può diventare un collo di
    bottiglia.
}

\fonteslide{37}

In pratica, però, ci sono pochissimi punti (se non nessuno) in un
programma medio dove questi problemi di efficienza sono determinanti; ma
in \textbf{sistemi real-time}, \textbf{routine matematiche core} o
\textbf{programmazione di sistema di basso livello} può essere il caso.

\fonteslide{38}

\nt{
    Questo è lo stesso trade-off discusso nella sezione trasversale del
    fascicolo di domande d'esame: C++ non-virtual-by-default segue lo
    \emph{zero-overhead principle}; Java virtual-by-default recupera il
    costo con la devirtualizzazione JIT (Class Hierarchy Analysis + inline
    cache).
}

%==============================================================
\chapter{L'esercizio del corso}\label{ch:esercizio}
%==============================================================

Le slide concludono con un esercizio completo su una gerarchia a sei
classi (\code{A}, \code{B}, \code{C}, \code{D}, \code{E}, \code{G}), che
mescola deliberatamente metodi virtual e non-virtual, override e hiding,
per costringere a ragionare riga per riga.

\fonteslide{40-45}

\begin{Example}{Le classi della gerarchia (estratto)}{}
\begin{lstlisting}[language=C++]
class A {
public:
    virtual void f() { prn("A::f()"); }
             void g() { prn("A::g()"); }
    virtual void h() { prn("A::h()"); }
             void i() { prn("A::i()"); }
};
class D : public A {
public:
    virtual void h() { prn("D::h()"); }
             void m() { prn("D::m()"); h(); }
    virtual void g() { prn("D::g()"); }
};
class G : public E {   // E deriva da D
public:
    virtual void f(int i) { prn("G::f(int)"); }
};
\end{lstlisting}
\end{Example}

\wc{Basta guardare quale metodo è virtual per sapere cosa stampa una chiamata}{
    No: bisogna anche controllare, per ogni chiamata, il \textbf{tipo
    statico} usato per farla (variabile, puntatore, riferimento). Tre
    righe di questo esercizio non compilano affatto, per \textbf{name
    hiding}: \code{G} dichiara \code{f(int)}, che nasconde ogni
    \code{f()} ereditato --- quindi \code{aG.f()} (senza argomenti) non
    trova nulla da chiamare, anche se una \code{f()} esiste più in alto
    nella gerarchia.
}

L'esercizio completo, da fare tu stesso prima di guardare la soluzione, è
in \code{prove/es\_00\_slide\_ABCDEG.cpp} (soluzione verificata in
\code{prove/soluzioni.md}).

%==============================================================
\chapter{Ereditarietà virtuale}\label{ch:vinherit}
%==============================================================

Le slide del corso non trattano l'ereditarietà virtuale: questo intero
capitolo si appoggia alla voce Wikipedia \emph{Virtual inheritance} e a
\emph{Thinking in C++} (Vol.~2, cap.~9, \emph{Multiple Inheritance}).

\fontewiki[Virtual inheritance]

\section{Il problema: sotto-oggetti duplicati}

Quando due classi \code{B} e \code{C} derivano da \code{A}, e una classe
\code{D} deriva da entrambe \code{B} e \code{C}, l'oggetto \code{D}
contiene \textbf{due copie} delle variabili membro di \code{A} --- una
via \code{B}, una via \code{C} --- accessibili indipendentemente tramite
scope resolution. Questo è il \textbf{diamond problem}.

\fontewiki[Virtual inheritance]
\fontelibro{cap.~9, ``Multiple Inheritance''}

\dfn{Ereditarietà virtuale}{
    Tecnica C++ che garantisce che \textbf{una sola copia} delle
    variabili membro di una classe base venga ereditata dalle classi
    derivate ``nipote''. Si dichiara con \code{virtual} nella lista delle
    basi: \code{class Derived : virtual public Base \{ ... \};}
}

\fontewiki[Virtual inheritance]
\fontelibro{cap.~9}

\section{Chi costruisce la base virtuale}

\mprop{La regola dell'inizializzazione}{
    La base virtuale viene inizializzata \textbf{sempre} dal costruttore
    della classe \textbf{most derived} --- mai da una classe intermedia,
    anche se quella intermedia scrive un mem-initializer per la base
    virtuale: se la classe intermedia non è quella davvero istanziata,
    quel mem-initializer viene \textbf{ignorato}. Questo impedisce che
    più classi intermedie inizializzino la base virtuale più volte.
}

\fontelibro{cap.~9, sezione ``Virtual base classes''}
\fonteprove{08\_virtual\_base\_ctor\_order.cpp}

\begin{Example}{Chi vince, tra i mem-initializer scritti a ogni livello}{}
\begin{lstlisting}[language=C++]
struct A { A(const char* who) { std::cout << who << " "; } };
struct B : virtual A { B() : A("da B") {} };
struct C : virtual A { C() : A("da C") {} };
struct D : B, C  { D() : A("da D"), B(), C() {} };
struct E : D     { E() : A("da E"), D()      {} };
int main() { E e; }
\end{lstlisting}
Output verificato: \code{A("da E") B() C() D() E()} --- solo il
mem-initializer di \code{E} (la classe most-derived) conta;
\code{"da B"}, \code{"da C"}, \code{"da D"} vengono tutti ignorati in
questa costruzione. Costruendo un \code{D} da solo, invece, vince
\code{"da D"}.
\end{Example}

\section{Dominanza}

\dfn{Dominanza}{
    Con basi virtuali esiste un \textbf{solo} sotto-oggetto condiviso: se
    solo una delle due classi intermedie fa l'override di un metodo
    virtuale della base condivisa, quell'override \textbf{domina}
    sull'implementazione originale, e diventa l'unico \emph{final
    overrider}: nessuna ambiguità.
}

\fonteprove{06\_diamond\_dominance.cpp}

\wc{Senza virtual, il diamante darebbe semplicemente lo stesso risultato ma più lento}{
    No: senza \code{virtual} nelle basi, esistono \textbf{due}
    sotto-oggetti \code{A} distinti, e la chiamata diventa
    \textbf{ambigua} --- errore di compilazione, non un semplice calo di
    prestazioni. Va risolta con una chiamata qualificata
    (\code{d.B::f();}).
}

\fonteprove{07\_diamond\_no\_virtual\_ambiguous.cpp}

Se \textbf{entrambe} le classi intermedie facessero l'override, non
esisterebbe più un final overrider unico: il programma diventa
\textbf{ill-formed alla definizione} della classe che le unisce (non al
punto di chiamata!) --- a meno che quella classe non fornisca il proprio
override, che allora domina su entrambi.

\fonteprove{es\_03\_diamond\_twist.cpp}

\section{Layout e costo in memoria}

\begin{Example}{L'esempio Animal/Mammal/WingedAnimal/Bat}{}
Con ereditarietà virtuale il layout di \code{Bat} diventa:
\begin{center}
\code{(vpointer, Mammal, vpointer, WingedAnimal, Bat, Animal)}
\end{center}
Il sotto-oggetto \code{Animal} finisce \textbf{in fondo} (non
all'inizio): \code{Mammal} e \code{WingedAnimal} devono conservare il
layout che avrebbero da soli, quindi l'unica copia condivisa di
\code{Animal} va in una posizione che non appartiene a nessuno dei due
rami. Servono \textbf{due vpointer}: ciascun ramo deve sapere ``quanto
dista da qui il mio \code{Animal}'' (\emph{vbase offset}), e questa
distanza non è nota a compile time (dipende da come il sotto-oggetto sarà
usato). La dimensione dell'oggetto cresce di due puntatori rispetto alla
versione non virtuale, in cambio dell'unicità di \code{Animal} e
dell'assenza di ambiguità.
\end{Example}

\fontewiki[Virtual inheritance]

\mprop{static\_cast è vietato, serve dynamic\_cast}{
    Un \code{static\_cast} da un puntatore a base virtuale verso la
    derivata è \textbf{vietato} (il compilatore dovrebbe sottrarre un
    offset che non conosce staticamente). Va usato \code{dynamic\_cast},
    che legge l'informazione dalla vtable a runtime --- l'unico caso in
    cui è \textbf{obbligatorio} pur non essendoci incertezza sul tipo
    effettivo.
}

\section{Metodi virtuali puri nel diamante}

Con ereditarietà virtuale, un metodo virtuale puro va implementato
\textbf{una sola volta} nella gerarchia: il ramo che non lo implementa
eredita l'implementazione dall'altra parte del diamante, per dominanza.

\fonteprove{09\_pure\_virtual\_diamond.cpp}

\section{Quando usarla}

\mprop{Il criterio}{
    L'ereditarietà virtuale serve quando l'ereditarietà esprime una
    \textbf{restrizione di un insieme} (un \code{Bat} \emph{è} un
    \code{Animal}, uno solo) e non una \textbf{composizione di parti}
    indipendenti (nel qual caso copie separate sono accettabili).
}

Esempio reale nella libreria standard: \code{basic\_istream} e
\code{basic\_ostream} derivano virtualmente da \code{basic\_ios}, così
\code{basic\_iostream} (da cui \code{stringstream}, \code{fstream}) ha un
solo stato di stream condiviso invece di due.

%==============================================================
\chapter{Oltre le slide: approfondimenti da Wikipedia}\label{ch:oltre}
%==============================================================

La voce Wikipedia \emph{Virtual function} elenca, in fondo alla pagina,
una serie di riferimenti specifici sul tema più insidioso di questo
capitolo: cosa succede chiamando una funzione virtuale dal costruttore o
dal distruttore. Questo capitolo li riprende.

\fontewiki

\section{Il corner case: una funzione virtuale pura può avere un corpo}

\wc{Una funzione virtuale pura (\code{= 0}) non può avere nessuna implementazione}{
    Falso, ed è un ``corner case'' citato esplicitamente da Wikipedia
    (Raymond Chen, \emph{``C++ corner case: you can implement pure
    virtual functions in the base class''}): \code{= 0} impone solo che
    le derivate concrete facciano l'override, non vieta un corpo. Quel
    corpo è chiamabile solo per chiamata qualificata --- o, come nel caso
    seguente, dal costruttore della classe stessa.
}

\fontewiki

\begin{Example}{Verificato: cosa succede DAVVERO chiamando una pura dal ctor}{}
Contrariamente a quanto ci si aspetterebbe (un crash a runtime), il
comportamento reale osservato (g++ 15.2.1, compilando esattamente il
codice sotto) è un \textbf{errore di link}, non un crash a runtime:
\begin{lstlisting}[language=C++]
struct Base {
    virtual void f() = 0;
    Base() { f(); }             // chiamata dal ctor
};
struct Derivata : Base {
    void f() override { std::cout << "Derivata::f\n"; }
};
// g++: undefined reference to `Base::f()'
\end{lstlisting}
Il motivo: dentro il costruttore di \code{Base}, il compilatore \emph{sa}
che il tipo dinamico dell'oggetto è esattamente \code{Base} (non ancora
\code{Derivata}): la chiamata \code{f()} viene \textbf{devirtualizzata a
compile time} in una chiamata diretta a \code{Base::f()}. Ma
\code{Base::f()} non esiste (pura, senza corpo): è il \textbf{linker},
non il runtime, ad accorgersene. Fornendo un corpo a \code{Base::f()}
(il corner case sopra), tutto funziona e quel corpo viene eseguito.
\end{Example}

\fonteprove{05\_pure\_virtual\_ctor.cpp}

\section{Il consiglio più citato: mai chiamare virtual da ctor/dtor}

Tra i riferimenti in fondo alla voce Wikipedia compare l'articolo di
Scott Meyers, \emph{``Never Call Virtual Functions during Construction or
Destruction''} (il capitolo 9 della terza edizione di \emph{Effective
C++}): la regola pratica --- indipendentemente dal fatto che il metodo sia
puro o abbia un corpo --- è di non farlo mai, perché il risultato
(dispatch sulla classe ``corrente'', non su quella finale) è quasi sempre
contro-intuitivo per chi legge il codice.

\fontewiki

\fonteprove{04\_ctor\_dtor\_virtual\_call.cpp}

\section{Argomenti collegati (See also)}

La voce Wikipedia elenca, tra gli argomenti collegati: \emph{Abstract
method}, \emph{Inheritance}, \emph{Superclass}, \emph{Virtual
inheritance} (capitolo~\ref{ch:vinherit}), \emph{Virtual class},
\emph{Interface}, \textbf{Component Object Model (COM)} e \emph{Virtual
method table} (capitolo~\ref{ch:vtable}).

\fontewiki

\nt{
    Il \textbf{Component Object Model} di Microsoft è l'esempio più noto
    di VFT usata come formato di interoperabilità binaria \emph{fuori}
    da un singolo linguaggio: un'interfaccia COM è, letteralmente, un
    puntatore a una VFT con un layout standardizzato --- lo stesso
    meccanismo di questo capitolo, usato per far comunicare linguaggi
    diversi (C++, VB, ecc.) tramite un contratto binario comune.
}

%==============================================================
\chapter{Overriding in Java}\label{ch:java-full}
%==============================================================

Questo capitolo si appoggia sulla voce Wikipedia \emph{Method overriding}
(sezione Java), sulla parte già vista dalle slide (capitolo~\ref{ch:javacpp})
e sul file verificato \code{prove/17\_JavaOverrideAnnotationAndCovariance.java}.

\fontewiki[Method overriding]

\section{Overriding automatico, e l'annotazione \code{@Override}}

\mprop{In Java l'overriding è automatico}{
    A differenza del C++, dove serve dichiarare esplicitamente
    \code{virtual}, in Java l'overriding avviene \textbf{automaticamente}
    ogni volta che un metodo di una sottoclasse combina lo stesso nome, gli
    stessi parametri e lo stesso return type (o uno covariante) di un
    metodo della superclasse.
}

\fontewiki[Method overriding]

L'annotazione \code{@Override} è \textbf{opzionale ma raccomandata}: non
cambia il comportamento a runtime, ma dice al compilatore ``questo metodo
deve essere un override di qualcosa: se non lo è, dammi errore''.

\fontewiki[Method overriding]

\begin{Example}{Perché \code{@Override} va sempre usata}{}
\begin{lstlisting}[language=Java]
class Animal { public void makeSound() { System.out.println("..."); } }
class Dog extends Animal {
    @Override
    public void makeSound() { System.out.println("bau"); }

    // Un typo nel nome del metodo, SENZA @Override, compila
    // silenziosamente: diventa un OVERLOAD inutile, non un override.
    // public void makeSund() { System.out.println("bau (typo)"); }
    //
    // Con @Override, lo stesso typo da' un errore di compilazione:
    // "method does not override a method from its superclass".
}
\end{lstlisting}
\end{Example}

\fonteprove{17\_JavaOverrideAnnotationAndCovariance.java}

\section{Cosa non può essere overridden}

\code{final}, \code{private} e \code{static} non possono essere
overridden; una classe \code{final} non può essere estesa affatto.

\fontewiki[Method overriding]
\fonteslide{16}

\wc{Un metodo private in una sottoclasse con lo stesso nome della base fa l'override}{
    No: i metodi \code{private} non sono ereditati né virtual
    (invocati con \code{invokespecial}, non con dispatch dinamico). Un
    metodo omonimo nella sottoclasse è un metodo del tutto
    \textbf{indipendente}, non un override --- esattamente come i metodi
    \code{private} in C++ (capitolo sul pattern NVI del fascicolo
    d'esame).
}

\section{super: chiamare l'implementazione della base}

Java usa la keyword \code{super} per invocare esplicitamente
l'implementazione della superclasse --- concettualmente equivalente al
\code{Base::f()} del C++ (che invece usa lo scope-resolution operator, non
una keyword dedicata).

\fontewiki[Method overriding]

\section{I campi NON sono polimorfici}

\wc{Se un metodo overridden dispatcha sul tipo reale, anche i campi si comportano cosi'}{
    No. I \textbf{campi} sono risolti a compile time sul tipo
    \textbf{statico} della variabile/riferimento, indipendentemente dal
    tipo reale dell'oggetto --- il polimorfismo in Java riguarda
    \emph{solo} i metodi.
}

\begin{lstlisting}[language=Java]
class A { String x = "A"; String get() { return x; } }
class B extends A { String x = "B"; String get() { return x; } }
A a = new B();
System.out.println(a.x);       // "A": campo, risolto sul tipo statico A
System.out.println(a.get());   // "B": metodo, risolto sul tipo dinamico B
\end{lstlisting}

\code{B.x} non sostituisce \code{A.x}: lo \textbf{nasconde} (nell'oggetto
esistono due campi distinti, accessibili con \code{super.x} o un cast).

\section{Return type covariante}

Da Java 5, un override può restituire un \textbf{sottotipo} del return
type dichiarato nella base (non necessariamente lo stesso tipo esatto).

\fonteprove{17\_JavaOverrideAnnotationAndCovariance.java}

\begin{lstlisting}[language=Java]
class Animal2 { Animal2 clone2() { return new Animal2(); } }
class Dog2 extends Animal2 {
    @Override
    Dog2 clone2() { return new Dog2(); }   // Dog2 e' un sottotipo di Animal2
}
\end{lstlisting}

\section{La chiamata a un metodo overridabile dal costruttore}

Il punto di confronto più importante con il C++ (capitolo~\ref{ch:oltre}):
in Java il tipo dinamico di un oggetto è \textbf{già quello finale} fin
dal primo istante della costruzione (a differenza del C++, dove ``cresce''
durante la costruzione, capitolo~\ref{ch:vtable}). Chiamare un metodo
overridabile dal costruttore della base dispatcha quindi
\textbf{correttamente} sulla sottoclasse --- ma i campi di quella
sottoclasse valgono ancora il \textbf{default} (anche se \code{final}!),
perché i suoi inizializzatori girano solo \emph{dopo} il ritorno da
\code{super()}.

\fonteprove{10\_VirtualFunctionsInJava.java}

\nt{
    Regola pratica identica a quella del C++ (Scott Meyers,
    capitolo~\ref{ch:oltre}), anche se il meccanismo che la causa è
    diverso: \textbf{mai chiamare, dal costruttore, un metodo che una
    sottoclasse potrebbe overridare}, in nessuno dei due linguaggi.
}

%==============================================================
\chapter{Overriding in Python}\label{ch:python-full}
%==============================================================

Questo capitolo si appoggia sulla voce Wikipedia \emph{Method overriding}
(sezione Python) e sui file verificati \code{prove/14}--\code{16}.

\fontewiki[Method overriding]

\section{Nessuna keyword, perché non serve}

In Python \textbf{ogni} metodo è dinamicamente dispatchato per
costruzione: il lookup di un attributo segue sempre l'MRO (\emph{Method
Resolution Order}) del tipo \textbf{reale} dell'oggetto. Non esiste una
versione ``non virtuale'' tra cui scegliere, quindi non c'è nulla da
dichiarare esplicitamente --- a differenza di C++ (dove serve
\code{virtual}) e, in misura minore, di Java (dove tutto è virtual di
default, ma esiste comunque \code{final} per vietarlo).

\fonteprove{11\_python\_no\_virtual\_keyword.py}

\section{super() e l'MRO: ereditarietà multipla cooperativa}

Wikipedia nota che Python usa \code{super(Subclass, self).method()} (o,
in Python 3, la forma semplificata \code{super().method()}) per invocare
l'implementazione della classe base.

\fontewiki[Method overriding]

\wc{super() chiama sempre ``la superclasse''}{
    Non è la superclasse: è il \textbf{prossimo elemento nell'MRO
    dell'oggetto}. In una gerarchia con ereditarietà multipla, questo è
    ciò che garantisce che un metodo comune a più rami del diamante venga
    eseguito \textbf{una sola volta}, non una per ramo.
}

\begin{Example}{Il diamante cooperativo, verificato}{}
\begin{lstlisting}[language=Python]
class A:
    def f(self): print("A.f")
class B(A):
    def f(self): print("B.f"); super().f()
class C(A):
    def f(self): print("C.f"); super().f()
class D(B, C):
    def f(self): print("D.f"); super().f()

print(D.__mro__)   # (D, B, C, A, object)
D().f()
\end{lstlisting}
Output verificato: \code{D.f B.f C.f A.f} (poi il ``dopo'' di ciascuno, in
ordine inverso) --- \code{A.f} viene eseguito \textbf{una sola volta}.
Dentro \code{B.f}, con \code{type(self) is D}, il successivo nell'MRO
dopo \code{B} è \code{C}, non \code{A}: è per questo che \code{super()} in
\code{B} non salta direttamente ad \code{A}.
\end{Example}

\fonteprove{14\_python\_mro\_c3.py}

\section{Name mangling: l'equivalente di un metodo non virtuale}

Python non ha \code{private} né un modo di impedire l'override in senso
stretto, ma il \textbf{doppio underscore} iniziale (\code{\_\_nome}) dà un
effetto simile tramite il \textbf{name mangling}: il compilatore riscrive
\code{self.\_\_helper} in \code{self.\_NomeClasse\_\_helper}, in base alla
classe \textbf{dentro cui} il codice è scritto (non alla classe reale
dell'oggetto a runtime).

\begin{Example}{Perché \code{B.\_\_helper} non fa l'override di \code{A.\_\_helper}}{}
\begin{lstlisting}[language=Python]
class A:
    def __init__(self): self.__helper()   # diventa self._A__helper()
    def __helper(self): print("A helper")
class B(A):
    def __helper(self): print("B helper")  # diventa _B__helper: nome diverso!
B()   # stampa "A helper", non "B helper"
\end{lstlisting}
I due metodi hanno nomi mangled \textbf{diversi} (\code{\_A\_\_helper} e
\code{\_B\_\_helper}): non c'è alcuna collisione, quindi nessun override.
Non è un meccanismo di sicurezza: il nome mangled resta accessibile da
fuori (\code{b.\_A\_\_helper()}).
\end{Example}

\fonteprove{15\_python\_name\_mangling.py}

\section{Metodi ``abstract'': @abstractmethod da solo non basta}

\wc{@abstractmethod impedisce di istanziare una classe che non fa l'override}{
    Da solo, no: \code{@abstractmethod} si limita a impostare
    \code{\_\_isabstractmethod\_\_ = True} sulla funzione. Il controllo
    che impedisce l'istanziazione è fatto dalla \textbf{metaclasse}
    \code{ABCMeta} (che popola \code{\_\_abstractmethods\_\_} alla
    creazione della classe): senza ereditare da \code{abc.ABC} (o
    impostare esplicitamente quella metaclasse), l'istanziazione riesce
    comunque, senza alcun errore.
}

\fonteprove{16\_python\_abc\_not\_enforced.py}

\section{Il costruttore e il dispatch}

Come in Java (capitolo~\ref{ch:java-full}), il tipo dinamico di
\code{self} è già quello finale dal primo istante della costruzione: una
chiamata a un metodo overridabile da \code{\_\_init\_\_} dispatcha
correttamente sulla sottoclasse. Vale però lo stesso avvertimento: se
quel metodo legge un attributo che l'\code{\_\_init\_\_} della sottoclasse
non ha ancora assegnato, il risultato può essere sbagliato (o, in Python,
un \code{AttributeError} se l'attributo non esiste affatto prima della
prima assegnazione).

\fonteprove{11\_python\_no\_virtual\_keyword.py}

%==============================================================
\chapter{Polimorfismo manuale in C}\label{ch:c-full}
%==============================================================

Il C non ha funzioni virtuali, classi, né alcuna forma di overriding nel
linguaggio: la voce Wikipedia \emph{Method overriding} non ha nemmeno una
sezione dedicata al C, per questo motivo. Questo capitolo traduce
\textbf{letteralmente a mano} il meccanismo VFT/VFTP del
capitolo~\ref{ch:vtable}, verificando col compilatore reale sia il
funzionamento sia i modi in cui si rompe.

\section{La traduzione letterale di VFT e VFTP}

\begin{Example}{Dispatch scritto a mano, verificato}{}
\begin{lstlisting}[language=C]
/* la VFT: una tabella di puntatori a funzione, UNA per classe */
typedef struct {
    void   (*draw)(const void *self);
    double (*area)(const void *self);
} ShapeVT;

/* la "classe base": porta solo il VFTP */
typedef struct { const ShapeVT *vptr; } Shape;

/* la "classe derivata": incorpora Shape come PRIMO membro */
typedef struct { Shape base; double r; } Circle;

static double circle_area(const void *self) {
    const Circle *c = (const Circle *)self;
    return 3.14159265 * c->r * c->r;
}
static const ShapeVT CIRCLE_VT = { circle_draw, circle_area };

Circle *circle_new(double r) {
    Circle *c = malloc(sizeof *c);
    c->base.vptr = &CIRCLE_VT;   /* il "costruttore" installa il vptr a mano */
    c->r = r;
    return c;
}
/* dispatch "virtuale": gli stessi 5 passi della teoria (cap. 3) */
double shape_area(const Shape *s) { return s->vptr->area(s); }
\end{lstlisting}
Un ciclo su un array di \code{Shape*} misti (\code{Circle}, \code{Rectangle})
chiama \code{shape\_area} \textbf{senza nessun \code{if}/\code{switch}
sul tipo}: esattamente il late binding della teoria, scritto a mano.
\end{Example}

\fonteprove{12\_manual\_vtable\_in\_c.c}

\mprop{Il primo membro deve essere la base}{
    Lo standard C garantisce l'assenza di padding all'inizio di una
    struct: quindi \code{(Shape*)circle == \&circle->base}, e la
    conversione derivata$\to$base è un no-op, senza bisogno di alcun
    aggiustamento del puntatore (nessun \emph{thunk}, a differenza
    dell'ereditarietà multipla C++).
}

\section{Chi inizializza il vptr: qui nessuno, se non tu}

\wc{Se dimentico di inizializzare il vptr, il programma comunque funziona un po'}{
    No: senza quella riga, il vptr è memoria non inizializzata
    (\code{malloc} non azzera la memoria). Chiamare una funzione
    attraverso un vptr indeterminato è undefined behavior, e nella
    pratica \textbf{crasha davvero} (verificato: \code{Segmentation
    fault}).
}

In C++, questa riga di inizializzazione il programmatore \textbf{non la
scrive mai a mano}: il compilatore la genera automaticamente in ogni
costruttore, per ogni livello della gerarchia (capitolo~\ref{ch:vtable}).
In C, dimenticarla è un errore reale e frequente.

\fonteprove{13\_c\_no\_forgotten\_vptr.c}

\section{Esempi reali: non è un esercizio accademico}

Questo identico schema (struct con un puntatore a una tabella di
puntatori a funzione) è usato da progetti C reali su larga scala:

\begin{itemize}
    \item il \textbf{VFS} (Virtual File System) del kernel Linux, dove
          \code{struct file\_operations} è, letteralmente, una VFT: ogni
          filesystem (ext4, procfs, un driver...) fornisce la propria
          tabella, e il kernel dispatcha tramite quella tabella senza
          sapere di che filesystem si tratti;
    \item \textbf{GObject} (la base di GTK): implementa un intero sistema
          di classi con ereditarietà single/multipla (via interfacce) e
          dispatch virtuale, tutto scritto a mano sopra il C puro,
          esattamente con questo schema.
\end{itemize}

\nt{
    Da distinguere dallo schema di \code{qsort(arr, n, sizeof(int), cmp)}:
    lì il puntatore a funzione (\code{cmp}) è passato \textbf{per
    parametro alla singola chiamata}, non associato al \emph{tipo} degli
    oggetti ordinati. È il pattern \emph{Strategy}, non l'overriding: non
    esiste alcuna ``classe'' che eredita quel comportamento.
}

%==============================================================
\chapter{Riferimenti e link aggiuntivi}\label{ch:riferimenti}
%==============================================================

Come indicato dal docente: Wikipedia va bene come riferimento, ma è
meglio eseguire gli esempi ed approfondire con i link che Wikipedia
stessa elenca in fondo alla pagina, più il libro \emph{Thinking in C++}.

\section{Dalla voce ``Virtual function''}

\textbf{See also}: Abstract method, Inheritance, Superclass, Virtual
inheritance, Virtual class, Interface, Component Object Model, Virtual
method table.

\textbf{Riferimenti specifici sul tema ctor/dtor} (il più utile per
questo capitolo): Raymond Chen, \emph{``What is \_\_purecall?''}; Scott
Meyers, \emph{``Never Call Virtual Functions during Construction or
Destruction''}; Raymond Chen, \emph{``C++ corner case: you can implement
pure virtual functions in the base class''}; S.G. Ganesh, \emph{``Joy of
Programming: Calling Virtual Functions from Constructors''}.

\section{Dalla voce ``Virtual inheritance''}

Milea, \emph{``Solving the Diamond Problem with Virtual Inheritance''}
(Cprogramming.com); McArdell, \emph{``C++/What is virtual
inheritance?''}.

\section{Dalla voce ``Method overriding''}

\textbf{See also}: Implementation inheritance, Inheritance semantics,
Method overloading, Polymorphism in object-oriented programming,
Template method pattern, Virtual inheritance.

\section{Thinking in C++ (Bruce Eckel) --- libero, due volumi}

\begin{itemize}
    \item Vol.~2, cap.~9 --- \emph{Multiple Inheritance}, sezione
          \emph{Virtual base classes}: \url{https://web.mit.edu/merolish/ticpp/TicV2.html\#_Toc53985825}
    \item Vol.~2, cap.~5 --- \emph{Templates in Depth} (per il capitolo
          sui template: non-type/template template parameter, SFINAE,
          CRTP, metaprogramming): \url{https://web.mit.edu/merolish/ticpp/TicV2.html\#_Toc53985706}
\end{itemize}

%==============================================================
\chapter{Fonti}\label{ch:fonti}
%==============================================================

\begin{itemize}
    \item Slide del corso: \emph{Virtual Functions and Late Bindings}
          (\code{20250922.ExtraLectureOnVF.pdf}), 45 pagine --- fonte
          primaria dei capitoli~\ref{ch:intro}--\ref{ch:esercizio}.
    \item Wikipedia, \emph{Virtual function},
          \url{https://en.wikipedia.org/wiki/Virtual_function} --- fonte
          primaria del capitolo~\ref{ch:oltre}.
    \item Wikipedia, \emph{Virtual inheritance},
          \url{https://en.wikipedia.org/wiki/Virtual_inheritance} ---
          fonte primaria del capitolo~\ref{ch:vinherit}.
    \item Wikipedia, \emph{Method overriding},
          \url{https://en.wikipedia.org/wiki/Method_overriding} ---
          fonte primaria dei capitoli~\ref{ch:java-full} e
          \ref{ch:python-full}.
    \item Bruce Eckel, \emph{Thinking in C++}, Vol.~2 (\emph{Practical
          Programming}), cap.~9 --- fonte di supporto per il
          capitolo~\ref{ch:vinherit}.
    \item Cartella \code{prove/} (questo stesso computer): ogni esempio
          citato con \fonteprove{} è stato compilato ed eseguito
          davvero, non solo descritto.
\end{itemize}

\end{document}
